\documentclass{beamer}

\def\Tiny{\fontsize{6pt}{6pt}\selectfont}
\def\supertiny{\fontsize{4pt}{4pt}\selectfont}

\mode<presentation>
{
  \usetheme{Warsaw}
  % \setbeamercovered{transparent}
  \usecolortheme{crane}
}

\usepackage{graphicx, ifthen, listings, fancyvrb}

\usepackage[czech]{babel}
% \usefonttheme{professionalfonts}
\usepackage{times}
\usepackage{amsmath}
\usepackage[utf8]{inputenc}
\usepackage{wrapfig}

\usepackage[T1]{fontenc}

\lstset{ basicstyle=\tiny, stringstyle=\ttfamily, showstringspaces=false }

\everymath{\displaystyle}

\setbeamerfont{frametitle}{size=\large}
\setbeamerfont{subsection in toc}{size=\scriptsize}

\makeatletter\newenvironment{blackbox}{%
   \begin{lrbox}{\@tempboxa}\begin{minipage}{0.95\columnwidth}}{\end{minipage}\end{lrbox}%
   \colorbox{black}{\usebox{\@tempboxa}}
}\makeatother

\title[IMF (8)]{Informatika pro moderní fyziky (8)\\Regulární výrazy. Tvorba textových dokumentů.}

\author[Franti\v{s}ek HAVL\r{U}J, ORF ÚJV Řež]{Franti\v{s}ek HAVL\r{U}J\\{\scriptsize \emph{e-mail: haf@ujv.cz}}}

\date{akademický rok 2019/2020\\13. listopadu 2019}

\institute[ORF ÚJV Řež]
{ÚJV Řež\\oddělení Reaktorové fyziky a podpory palivového cyklu}

\AtBeginSection[]
{
\begin{frame}<beamer>
\frametitle{Obsah}
\tableofcontents[currentsection,hideothersubsections]
\end{frame}
}

\begin{document}

\begin{frame}
  \titlepage
\end{frame}

\begin{frame}
  \tableofcontents
\end{frame}

\section{Co jsme se naučili minule}

\begin{frame}{}
  \begin{itemize}
    \item tabulky v latexu
    \item základy erb
  \end{itemize}
\end{frame}

\section{Najdi to, nevím co: regex}

\begin{frame}{Úkoly}
  \begin{itemize}
    \item najít všechna celá čísla: jaké je nejmenší a největší číslo a kolik jich je
    \item najít všechny emailové adresy: vypište je do souboru seskupené podle domény
    \item vlnka: nedělitelné mezery v textu před k, s, v, z (výsledek vypište do souboru)
    \item nahradit desetinné tečky čárkami (výsledek vypište do souboru)
  \end{itemize}
\end{frame}

\begin{frame}{Jak na to: regulární výrazy (regexy)}
  \begin{itemize}
    \item můžu definovat chytrou `masku'
    \item nejjednodušší: /text/ - text
    \item skupiny písmen: /1[abc]/ - 1a, 1b ...
    \item speciální skupiny: /\\da/ - 1a, 2a, 3a ...
    \item opakování: +, *
    \item vytažení části výrazu do tzv. capture group -- stačí uzávorkovat
    \item hračka a vysvětlení: rubular.com
  \end{itemize}
\end{frame}

\begin{frame}{Regexy v Ruby}
  \begin{itemize}
    \item \texttt{String\#match} vrací MatchData objekt, m[0] je ten samotný řetězec, m[i] jsou skupiny
    \item \texttt{String\#gsub} nahrazuje; můžu použít s blokem, nicméně tam nemám capture groups (ale můžu použít \texttt{Regexp.last\_match})
    \item pokud použiju formu bez bloku, mám capture groups v \textbackslash \textbackslash 1 a \textbackslash \textbackslash 2
    \item \texttt{String\#scan} vrací pole všech výskytů
  \end{itemize}
\end{frame}


\section{Výroba dokumentu v~praxi - ERb}

\begin{frame}{Úkoly na dnešek}
  \begin{itemize}
    \item násobilku 10x10 z minula přepracovat tak, aby se vygenerovala 7x7, 9x9, 11x11 - použití parametrů v ERb
    \item přepracovat náš krásný report z JE Třeskoprsky s použitím ERb 
    \item tabulky nicméně zajisté předělávat nebudeme
  \end{itemize}
\end{frame}

\begin{frame}[fragile]{Základní syntaxe ERb (1)}
  \begin{block}{ }
    Jakýkoli Ruby příkaz, přiřazení, výpočet ...
    \scriptsize
    \begin{verbatim}
      <% a = b + 5 %>
      <% list = ary * ", " %>
    \end{verbatim}
  \end{block}
\end{frame}

\begin{frame}[fragile]{Základní syntaxe ERb (2)}
  \begin{block}{ }
    Pokud chci něco vložit, stačí přidat rovnítko
    \scriptsize
    \begin{verbatim}
      <%= a %>
      <%= ary[1] %>
      <%= b + 5 %>
    \end{verbatim}
  \end{block}
\end{frame}

\begin{frame}[fragile]{Základní syntaxe ERb (3)}
  \begin{block}{ }
    Radost je možnost použít bloky a tedy i iterátory apod. v~propojení s~vkládaným textem:
    \scriptsize
    \begin{verbatim}
      <% (1..5).each do |i| %>
      Number <%= i %>
      <% end %>
      <% ary.each do |x| %>
      Array contains <%= x %>
      <% end %>
    \end{verbatim}
  \end{block}
\end{frame}

\begin{frame}{Co je podstatné}
  \begin{itemize}
    \item nezapomenout, kdy jsem v Ruby, kdy v ERb, kdy v latexu a vždy dodržet zadanou syntaxi
    \item mít na paměti, jak dostat data ze skriptu do šablony
    \item potřebuju mít načtený seznam všech kampaní do nějakého pole a to poslat do šablony
  \end{itemize}
\end{frame}

\begin{frame}{A to je vše, přátelé!}
  \begin{center}
    \includegraphics[width=0.8\textwidth]{looney_tunes}
  \end{center}
\end{frame}

\end{document}
